% 小抄示例（微观经济学，只能带一张A4）。编译：python scripts/compile_tex.py sheet_example.tex --max-pages 1
\documentclass[a4paper,fontset=windows]{ctexart}
\ifdefined\SBSIZE\else\def\SBSIZE{9}\fi
\ifdefined\SBLEAD\else\def\SBLEAD{10.8}\fi
\ifdefined\SBMARGIN\else\def\SBMARGIN{10}\fi
\usepackage[margin=\SBMARGIN mm]{geometry}
\usepackage{amsmath,amssymb,enumitem,multicol,booktabs,xcolor}
\makeatletter
\renewcommand\normalsize{\@setfontsize\normalsize{\SBSIZE}{\SBLEAD}%
  \abovedisplayskip 2\p@ \@plus1\p@ \@minus1\p@
  \belowdisplayskip \abovedisplayskip
  \abovedisplayshortskip 1\p@ \belowdisplayshortskip 1\p@
  \let\@listi\@listI}
\makeatother
\normalsize
\pagestyle{empty}
\setlength{\parindent}{0pt}
\setlength{\parskip}{1pt}
\setlength{\columnsep}{5pt}
\setlength{\columnseprule}{0.2pt}
\setlength{\multicolsep}{2pt}
\setlist{nosep,leftmargin=1.1em,label=\textbullet}
\newlength{\sbw}
\newcommand{\kw}[1]{\ifdefined\SBFILL\settowidth{\sbw}{#1}\underline{\hspace*{\dimexpr\sbw+0.6em\relax}}\else\textbf{#1}\fi}
\newcommand{\hd}[1]{\par\vspace{2pt}\textbf{【#1】}}

\begin{document}
{\bfseries 微观经济学　2027-01-08}\hfill{\color{gray}A4一页}
\begin{multicols}{3}
\hd{需求与供给}
\begin{itemize}
\item \kw{需求定律}：价格升，需求量降；\kw{供给定律}：价格升，供给量升
\item 均衡：$Q_d(P^*)=Q_s(P^*)$。价格高于均衡 \kw{过剩}，低于均衡 \kw{短缺}
\item 价格变 $\to$ 沿曲线移动；收入、偏好、替代品价格变 $\to$ 曲线移动
\item 需求右移：$P\uparrow$ $Q\uparrow$；供给右移：$P\downarrow$ $Q\uparrow$
\item 题型：令 $Q_d=Q_s$ 解 $P^*$，代回求 $Q^*$
\end{itemize}

\hd{价格弹性}
\begin{itemize}
\item $E_d=\left|\frac{\Delta Q/Q}{\Delta P/P}\right|$，$>1$ \kw{富有弹性}，$<1$ \kw{缺乏弹性}，$=1$ 单位弹性
\item $TR=P\times Q$。涨价：富有弹性 $TR\downarrow$，缺乏弹性 $TR\uparrow$
\item 因素：\kw{替代品}、占收入比重、\kw{时间}（越长越有弹性）
\item 例：涨价20\%、销量降30\%，$E_d=1.5$，富有弹性，$TR$ 降
\item 往年题每年一道计算题
\end{itemize}

\hd{消费者剩余}
\begin{itemize}
\item 愿意支付的最高价$-$实际价格；图上是需求曲线以下、价格线以上的面积
\item 线性需求：$CS=\frac12(P_{\max}-P)Q$
\item 例：$P=20-Q$，均衡价8，$Q=12$，$CS=\frac12\times12\times12=72$
\item 价格降，$CS$ 增；总剩余$=CS+PS$
\end{itemize}
\end{multicols}
\end{document}
